\documentclass[compacto,10pt]{pubciencias-notas}

% -- Paquetes adicionales
\usepackage{comandosEPN}

% -- Datos del libro
\autor{Departamento de Formación Básica}
\materia{Álgebra Lineal}
\nota{Semana 06 Clase 02}
\fecha{Semestre 2025-A}

\logouno[1.1cm]{Logo_EPN.eps}
\logodos[1.7cm]{Logo_DFB.eps}

\definecolor{colordef}{cmyk}{0.81,0.62,0.00,0.22}

\begin{document}

\encabezado

\vspace*{-8mm}

\setcounter{section}{2}
\section{Determinantes}
\setcounter{subsection}{1}
\subsection{Resolución de sistemas de ecuaciones}

\begin{teo}[Equivalencias no singulares]
Sea $A\in\Mat[\K]{n}{n}$, los siguientes enunciados son equivalentes.
\begin{enumerate}
\item $A$ es no singular;
\item el sistema $Ax=0$ tiene solamente la solución trivial;
\item $A$ es equivalente por filas a $I_{n}$;
\item $\rang(A)=n$;
\item el sistema lineal $Ax=b$ tiene una solución única para cada $b\in\Mat[\K]{n}{1}$; y
\item $\det(A)\neq0$.
\end{enumerate}
\end{teo}
\begin{ejer}
    Sea $\alpha\in\mathbb{R}$. Resuelva el sistema de ecuaciones:
     \[
        \l\{
        \begin{array}{c@{\,}c@{\,}c@{\,}c@{\,}c@{\,}c@{\,}c}
        \alpha x_1  & + & 2x_2 &= & 0 \\
        2x_1  & + & \alpha x_2 &= & 0
        \end{array}
        \r.
    \]
    ¿Para qué valores de $\alpha$ el sistema tiene solución trivial?
\end{ejer}
\begin{ejer}
    Sea $\alpha\in\mathbb{R}$. Dada la matriz
     \[
        A=\begin{pmatrix}
            1&2&3\\
            \alpha&1&0\\
            -1&2&1
        \end{pmatrix}
    \]
    ¿Para qué valores de $\alpha$ la matriz $A$ es no singular?
\end{ejer}
\subsection{Regla de Cramer}
\begin{teo}[Regla de Cramer]
    Sean $A\in\Mat[\K]{n}{n}$ y $b\in\Mat[\K]{n}{1}$. Consideremos el sistema de ecuaciones
    \[
        \l\{
        \begin{array}{c@{\,}c@{\,}c@{\,}c@{\,}c@{\,}c@{\,}c}
        a_{11} x_1 & + & a_{12} x_2 & + \, \cdots \, + & a_{1n} x_n & = & b_1, \\
        a_{21} x_1 & + & a_{22} x_2 & + \, \cdots \, + & a_{2n} x_n & = & b_2, \\
        \vdots &    &     \vdots & & \vdots & &  \vdots \\
        a_{n1} x_1 & + & a_{n2} x_2 & + \, \cdots \, + & a_{nn} x_n & = & b_n, \\
        \end{array}
        \r.
    \]
    y denotemos por $A$ a su matriz de coeficientes y $b$ a su matriz de constantes. Si $\det(A) \neq 0$, el sistema tiene solución única y está dada por:
    \[
        x_j = \dfrac{\det(A_j)}{\det(A)},
    \]
    para $j\in I$ y donde $A_j$ es la matriz que se obtiene reemplazando la $j$-ésima columna de $A$ por $b$.
\end{teo}

\begin{ejem}
    Consideremos el sistema de ecuaciones:
    \[
        \l\{
        \begin{array}{c@{\,}c@{\,}c@{\,}c@{\,}c@{\,}c@{\,}c}
        2x_1 & + & x_2  & - & x_3  & = & 3 \\
        x_1  & + & x_2  & + & x_3  & = & 1 \\
        x_1  & - & 2x_2 & - & 3x_3 & = & 4 \\
        \end{array}
        \r.
    \]
    En donde:
    \[
        A = 
        \begin{pmatrix}
            2 & 1 & -1 \\
            1 & 1 & 1 \\
            1 & -2 & -3
        \end{pmatrix} \qquad y \qquad b =
        \begin{pmatrix}
            3 \\
            1 \\
            4
        \end{pmatrix}
    \]
    Se sigue que $\det(A) = 5$. Puesto que $\det(A) \not = 0$, el sistema tiene solución única, la cual está determinada por:
    \[
        x_1 = \dfrac{
        \begin{vmatrix}
            3 & 1 & -1 \\
            1 & 1 & 1 \\
            4 & -2 & -3
        \end{vmatrix}}{5} = 2
    \]
    \[
        x_2 = \dfrac{
        \begin{vmatrix}
            2 & 3 & -1 \\
            1 & 1 & 1 \\
            1 & 4 & -3
        \end{vmatrix}}{5} = -1
    \]
    \[
        x_3 = \dfrac{
        \begin{vmatrix}
            2 & 1 & 3 \\
            1 & 1 & 1 \\
            1 & -2 & 4
        \end{vmatrix}}{5} = 0
    \]
    Así, la solución única del sistema es $x = \begin{pmatrix}2&-1&0\end{pmatrix}^{\intercal}$.
\end{ejem}

\begin{ejer}
    Resuelva el sistema de ecuaciones:
 %   \[
  %      \l\{
   %     \begin{array}{c@{\,}c@{\,}c@{\,}c@{\,}c}
    %    -ax_1  & + & 2bx_2 & = & c \\
     %   2ax_1 & - & 5bx_2 & = & 3c
      %  \end{array}
     %   \r.
    %\]
    %donde $a$, $b$ y $c$ son números reales, con $a b \not = 0$

     \[
        \l\{
        \begin{array}{c@{\,}c@{\,}c@{\,}c@{\,}c@{\,}c@{\,}c}
        3x  & + & 2y & + & z &= & 2 \\
        2x  & + & 2y & + & z &= & 2\\
        x  & + & y & + & z &= & 1
        \end{array}
        \r.
    \]


\end{ejer}
\end{document}